We now introduce the important notion of \emph{conjugate elements}.  Two 
elements $A$ and $B$ are said to be conjugate to one another if
\[
A=CBC^{-1},
\]
where $C$ is also an element of the group.  Multiplying this equality by $C$ 
on the right and by $C^{-1}$ on the left gives the inverse relation 
$B=C^{-1}AC$.  An essential property of conjugacy is transitivity: if $A$ is 
conjugate to $B$, and $B$ to $C$, then $A$ and $C$ are conjugate as well.  In 
fact, from $B=P^{-1}AP$ and $C=Q^{-1}BQ$, with $P,Q$ elements of the group, it 
follows that $C=(PQ)^{-1}A(PQ)$.  One may therefore consider sets of elements 
that are mutually conjugate.  Such sets are called \emph{conjugacy classes}, 
or simply classes of the group.  Any one of its elements $A$ determines the 
entire class: once $A$ is given, form the products $GAG^{-1}$ as $G$ ranges 
over all elements of the group; the same member of the class may, of course, 
be obtained more than once.  In this way the whole group can be partitioned 
into classes, and each element evidently belongs to only one class.
\SourcePage{437}{440}
The identity element forms a class by itself, since $GEG^{-1}=E$ for every 
element $G$ of the group.  In an Abelian group the same is true of each 
element.  All elements of such a group commute by definition, so every 
element is conjugate only to itself and hence constitutes its own class.  It 
should be emphasized that a class other than $E$ is by no means a subgroup: 
this already follows from the fact that it does not contain the identity 
element.
